\section{Objetivos}
\subsection{Objetivo general}

Desarrollar una aplicación web distribuida para la gestión del refuerzo académico en la Universidad Técnica Estatal de Quevedo, mediante la integración de servicios independientes que intercambien información académica y gestionen las solicitudes, para facilitar la coordinación y el seguimiento de las sesiones entre estudiantes y docentes.

\subsection{Objetivos específicos}

\begin{enumerate}
\item Analizar el proceso de refuerzo académico para identificar las necesidades de sus participantes y delimitar las funciones y los datos que deberán coordinarse entre los servicios de la aplicación.
\item Diseñar la arquitectura distribuida de la aplicación, definiendo las responsabilidades de los servicios, los contratos de comunicación y el control de acceso a la información.
\item Implementar los servicios de la aplicación para gestionar el ciclo de atención del refuerzo académico y proporcionar la información académica necesaria para su funcionamiento.
\item Integrar los servicios mediante comunicación por red, contemplando el intercambio de datos, el tratamiento de fallos y las notificaciones a los participantes.
\item Evaluar el funcionamiento integrado, el rendimiento y la facilidad de uso de la aplicación distribuida mediante pruebas técnicas y tareas realizadas por usuarios representativos.
\end{enumerate}

\subsection{Indicadores de cumplimiento}

La evaluación final del proyecto utilizará los siguientes indicadores para comprobar el cumplimiento de cada objetivo específico.

\begin{longtable}{@{}p{0.10\textwidth}>{\raggedright\arraybackslash}p{0.80\textwidth}@{}}
\toprule
\textbf{Objetivo} & \textbf{Indicador y comprobación prevista}\\
\midrule
\endhead
OE1 & Se comprobará que el análisis cubra a los cinco perfiles: estudiante, docente, coordinación, gestor de espacios y administración. Para cada perfil se identificarán sus actividades y necesidades; cada grupo de datos utilizado se asociará con la función que lo administra.\\[0.15cm]
OE2 & El 100\,\% de los servicios del diseño tendrá una responsabilidad definida y cada interacción especificará emisor, receptor, datos, permisos y tratamiento de fallos. Se comprobará la correspondencia entre los diagramas y los contratos.\\[0.15cm]
OE3 & Se ejecutarán al menos dos servicios de forma independiente. Se probará el ciclo completo de una solicitud individual y otra grupal, con consulta académica, programación, reserva presencial, reprogramación, asistencia y cierre, incluyendo la variante virtual.\\[0.15cm]
OE4 & Se comprobarán cinco escenarios de consulta: dato válido, entrada inválida, dato inexistente, servicio académico no disponible y recuperación del servicio. Para cada cambio de estado se verificará el envío de la notificación a los destinatarios correspondientes.\\[0.15cm]
OE5 & Se medirán los tiempos de las operaciones principales con hasta 100 usuarios concurrentes y 10\,000 solicitudes, con una meta máxima de 5 segundos. Se observará la ejecución de tareas sin asistencia y se aplicará SUS, con una meta promedio de al menos 70 puntos.\\
\bottomrule
\end{longtable}
